%-------------------------
% Cover Letter - Graviton | Quantitative Research Intern (2028 Graduates)
% Gurugram, Haryana | 8-week internship
% Apply with: Cover Letter Collection/ojas_srivastava_graviton_quantitative_research_intern_2028_cover_letter.pdf
% Ojas Srivastava
%-------------------------
\documentclass[a4paper,11pt]{article}
\usepackage[empty]{fullpage}
\usepackage{geometry}
\usepackage[hidelinks]{hyperref}
\usepackage{parskip}
\usepackage[T1]{fontenc}

\geometry{left=2.5cm, top=2cm, right=2.5cm, bottom=2cm}
\pagestyle{empty}
\setlength{\parindent}{0pt}
\setlength{\parskip}{10pt}

\begin{document}

\begin{flushright}
\textbf{Ojas Srivastava} \\
+91-7424978046 \\
\href{mailto:srivastavaojas454@gmail.com}{srivastavaojas454@gmail.com} \\
\href{https://github.com/Ojas-Srivastava05}{github.com/Ojas-Srivastava05} \textbar{}
\href{https://linkedin.com/in/ojas-srivastava05}{LinkedIn} \\
Surat / Mumbai \textbar{} Available for on-site internship in Gurugram \\
\today
\end{flushright}

\vspace{4mm}
\textbf{Graviton Research Capital --- Quantitative Research Hiring} \\
Quantitative Research Intern (2028 Graduates) \\
Gurugram, Haryana, India

\vspace{2mm}
\textbf{Re: Application for Quantitative Research Intern --- 2028 Graduates}

Dear Hiring Team,

I am applying for the Quantitative Research Intern role for 2028 graduates in Gurugram. I am a penultimate-year B.Tech Artificial Intelligence student at SVNIT Surat (CGPA 9.20/10, graduating May 2028), with coursework in probability and statistics, linear algebra, and machine learning, and strong programming skills in Python and C++. What draws me to Graviton is the idea that trading edges begin as research questions---curiosity, statistical rigor, and evidence over instinct---and that an eight-week internship can mean owning work that is evaluated against real markets, not observed from the sidelines.

My strongest research-style project is LogiFlow (Google Solution Challenge 2026 Global Top 106). I treated delay prediction as an open-ended research problem: formulated the hypothesis, engineered features from timetable, route, and history on 15,650 train-day records, and built Gradient Boosting models evaluated with cross-validation (MAE 22.7 minutes; 81\% of predictions within 30 minutes). I iterated on experimental design, built a 100/100 automated validation suite, and shipped analytical tooling for ranking and explainability---practice moving from question to model to measured outcome. Separately, I have invested in Indian equities and mutual funds for over three years and write Python tooling for portfolio logging, attribution, and screening; that habit of testing ideas against data is how I prefer to think about markets.

I am comfortable working with large structured datasets, writing clean Python for analysis and modeling, and communicating technical results clearly (ACM SVNIT executive; Nexus mentor). Competitive programming (LeetCode Knight, rating 2048; Codeforces 1419) keeps my problem-solving sharp when the first approach is wrong. I would be excited to spend eight weeks at Graviton contributing to predictive research, learning from quantitative researchers and technologists, and holding myself to the standard that success means a complete project from hypothesis through evaluation.

Thank you for your consideration. I would welcome the opportunity to discuss how I can contribute.

Sincerely, \\
\vspace{8mm}
\textbf{Ojas Srivastava}

\end{document}
